\documentclass[12pt]{article}
\usepackage[utf8]{inputenc}
\usepackage[T1]{fontenc}
\usepackage{amsmath,amssymb}
\usepackage{geometry}
\geometry{margin=2.5cm}

\begin{document}

\noindent $\tau \overset{\text{def}}{=} (i,j)$

\noindent $\pi \overset{\text{def}}{=} \sigma\circ\tau$

\[
\pi(j) = \sigma\big(\tau(j)\big) = \sigma(i) = j
\]

\bigskip
\noindent If $\sigma(k)=k \Rightarrow \pi(k)=k \Rightarrow k\neq i, k\neq j \ \Rightarrow d_\pi<d_\sigma.$

\[
\pi(k) = \sigma\big(\tau(k)\big) = \sigma(k) = k.
\]

\noindent $\pi = \sigma\circ\tau \ \big| \ \tau.$

\bigskip
\noindent $\pi = \tau_1\circ\cdots\circ\tau_s$ \ (induction hyp.) \ $\big|\tau \ \Rightarrow \ \tau_1\circ\cdots\circ\tau_s\circ\tau=\sigma.$

\bigskip
\noindent If $\sigma\in S_n$ \ inv$(\sigma) \overset{\text{not}}{=} \text{card}\{(i,j) \mid 1\leq i<j\leq n,\ \sigma(i)>\sigma(j)\}.$

\noindent (= the number of inversions of $\sigma$.)

\[
\sigma = \begin{pmatrix}1&2&\cdots&i&=&j&\cdots&n\\ \sigma(1)&\sigma(2)&\cdots&\sigma(i)&\cdots&\sigma(j)&\cdots&\sigma(n)\end{pmatrix} \quad C_n^2.
\]

\noindent inv$(e)=0$.

\bigskip
\noindent\fbox{Def}: \ The number $\varepsilon(\sigma) = (-1)^{\text{inv}(\sigma)}$ is called the \underline{signature}.

\bigskip
\noindent $\sigma$ is even if $\varepsilon(\sigma)=1$.

\noindent \phantom{$\sigma$ is }odd if $\varepsilon(\sigma)=-1$.

\bigskip
\noindent\fbox{T} \ If $\sigma\in S_n$ and $\sigma=\tau_1\circ\tau_2\circ\cdots\circ\tau_m \Rightarrow$ the number $m$ and inv$(\sigma)$ have the same

\noindent parity.

\bigskip
\noindent\fbox{Lemma} \ (a) \ Let $\tau,\sigma\in S_n$, $\tau$ a transposition

\noindent Then inv$(\sigma)$ and inv$(\sigma\circ\tau)$ have different parities.

\noindent (b) Every transposition is an odd permutation.

\bigskip
\noindent\fbox{Proof of Lemma} \ $\tau=(i,j)$ \quad $1\leq i<j\leq n$ and $r=j-i\in\mathbb{N}^*$.

\bigskip
\noindent If $r=1$

\[
\sigma = \begin{pmatrix}1&\cdots&i&i+1&\cdots&n\\ \sigma(1)&\cdots&\sigma(i)&\sigma(i+1)&\cdots&\sigma(n)\end{pmatrix}
= \begin{pmatrix}1&\cdots&i&j&\cdots&n\\ \sigma(1)&\cdots&\sigma(i)&\sigma(j)&\cdots&\sigma(n)\end{pmatrix}
\]

\[
\pi=\sigma\circ\tau = \begin{pmatrix}1&\cdots&i&j&\cdots&n\\ \sigma(1)&\cdots&\sigma(j)&\sigma(i)&\cdots&\sigma(n)\end{pmatrix}
= \begin{pmatrix}1&\cdots&i&i+1&\cdots&n\\ \sigma(1)&\cdots&\sigma(i+1)&\sigma(i)&\cdots&\sigma(n)\end{pmatrix}
\]

\bigskip
\noindent $r=1 \Rightarrow$ inv$(\sigma\circ\tau) =
\begin{cases}
\text{inv}(\sigma)+1, & \text{if } \sigma(i)<\sigma(i+1)\\
\text{inv}(\sigma)-1, & \text{if } \sigma(i)>\sigma(i+1)
\end{cases}$

\bigskip
\noindent\underline{General case}: \ $\tau=(i,j) \Rightarrow j-i=r$, $r>1$.

\noindent $\pi$ is obtained from $\sigma$ by exchanging $\sigma(i)$ with $\sigma(j)$, achieved through a sequence of transpositions of adjacent positions.\footnote{this paragraph of the proof is very compressed in the original; the exact step (``$d_\pi$ ... exchange ...'') could not be reconstructed with full certainty, but the standard argument is the one given here.}

\[
2r-1 \text{ steps}
\]

\noindent each step changes the parity of inv (the case $r=1$) $\Rightarrow$ after $2r-1$ steps (an odd number)

\noindent the parity of inv$(\sigma)$ changes. \hfill different parities \quad Q.E.D.

\end{document}
